\section{Connexion nonadique, extension du calendrier et mémoire de transport}
\label{cap:calendario}

La chronologie du TPK combine une phase observable, un registre dodécaphasique
et une mémoire des transitions effectuées. Leur relation est déterminée par
la composition des déplacements. Le franchissement de la frontière de phase
incrémente un quotient ; la conservation de ce quotient distingue les
retours successifs. Cette structure fournit la base arithmétique de
l’holonomie nonadique et de sa représentation bilatérale.

Il convient de fixer dès le départ les trois objets qui interviennent. La phase
appartient à un cycle de neuf positions. Le calendrier fini retient douze
tours de ce cycle. L’histoire enrichie conserve en outre le nombre
entier de tours et les registres de position, d’orientation, de feuillet, de retenue et de
frontière. Chaque réduction conserve une observation déterminée de la même
évolution ; l’information disponible est précisée au moyen de ses projections.

\subsection{Incidence des modes arithmétiques}
\label{sec:cal-incidencia}

\subsubsection{Sélection, permanence et succession cyclique}

Soient $\Sigma$ et $\Pi$ les modes correspondant aux évaluations additive
et multiplicative d’APP. Le sélecteur de phase parcourt, dans chaque bloc de trois
événements, la succession $(+1,-1,0)$. Son action sur le mode $m$ est
\[
 +1:m\longmapsto\Sigma,\qquad
 -1:m\longmapsto\Pi,\qquad
 0:m\longmapsto m.
\]
L’événement neutre conserve le mode précédent. En partant de la clôture cyclique
du bloc, les trois modes enregistrés sont donc
\[
 (\Sigma,\Pi,\Pi).
\]
La troisième position conserve l’étiquette de mode multiplicatif
sélectionnée lors de l’événement précédent. Dans l’émetteur, cet événement neutre
incrémente uniquement le compteur $Z$. Cette distinction permet de
construire l’incidence directement à partir des transitions du calendrier.

\begin{proposicion}[Incidence modale du bloc tritique]
\label{prop:cal-incidencia}
En comptant les paires consécutives du bloc, y compris sa clôture, et en ordonnant
les modes comme $(\Sigma,\Pi)$, on obtient
\begin{equation}
 F_{\rm av}=\begin{pmatrix}0&1\\1&1\end{pmatrix}.
 \label{eq:cal-incidencia}
\end{equation}
Dans une succession de $3q$ événements composée de $q$ répétitions du bloc,
le décompte des transitions est $qF_{\rm av}$.
\end{proposicion}
\begin{proof}
Les paires ordonnées sont $\Sigma\to\Pi$, $\Pi\to\Pi$ et
$\Pi\to\Sigma$. Chacune apparaît une fois ; $\Sigma\to\Sigma$
apparaît zéro fois. Ce sont respectivement les entrées supérieure
droite, inférieure droite, inférieure gauche et supérieure gauche de la
matrice. Chaque répétition ajoute exactement les mêmes trois paires, y compris
le lien du dernier mode d’un bloc au premier du suivant.
La somme des décomptes est $qF_{\rm av}$.
\end{proof}

Les longueurs $9$, $27$ et $108$ produisent ainsi $3F_{\rm av}$,
$9F_{\rm av}$ et $36F_{\rm av}$. Ces matrices comptent les événements d’un
itinéraire périodique déjà fixé. Les puissances de $F_{\rm av}$ comptent
les concaténations admises par son incidence. Il s’agit de deux opérations
mathématiques distinctes sur la même matrice : somme d’occurrences et
composition de transitions.

\subsubsection{Concaténation des incidences et récurrence induite}

Définissons les entiers $a_0=0$, $a_1=1$ et
$a_{n+1}=a_n+a_{n-1}$ pour $n\geq1$. Leur apparition découle de la
multiplication de la matrice déjà construite.

\begin{proposicion}[Puissances de l’incidence modale]
\label{prop:cal-potencias}
Pour tout $n\geq1$,
\[
 F_{\rm av}^{\,n}=
 \begin{pmatrix}a_{n-1}&a_n\\a_n&a_{n+1}\end{pmatrix},
 \qquad
 a_{n-1}a_{n+1}-a_n^2=(-1)^n.
\]
\end{proposicion}
\begin{proof}
Pour $n=1$, la formule reproduit \eqref{eq:cal-incidencia} parce que
$a_2=1$. Si la formule vaut pour $n$, la multiplication par $F_{\rm av}$
donne
\[
 \begin{pmatrix}
 a_n&a_{n-1}+a_n\\a_{n+1}&a_n+a_{n+1}
 \end{pmatrix}
 =\begin{pmatrix}a_n&a_{n+1}\\a_{n+1}&a_{n+2}\end{pmatrix}.
\]
La récurrence prouve la première identité. En prenant les déterminants et en utilisant
$\det F_{\rm av}=-1$, on obtient la seconde.
\end{proof}

En particulier, $F_{\rm av}^{2}=F_{\rm av}+I$. Cette identité permet de
normaliser un générateur hyperbolique au moyen de
\begin{equation}
 J_{-1}=\frac{2F_{\rm av}^{2}-3I}{\sqrt5}
       =\frac{2F_{\rm av}-I}{\sqrt5},
 \qquad J_{-1}^{2}=I.
 \label{eq:cal-generador-hiperbolico}
\end{equation}
En effet, $(2F_{\rm av}-I)^2=4(F_{\rm av}+I)-4F_{\rm av}+I=5I$.
Le coefficient $5$ provient ici de l’identité matricielle. La représentation
quadratique et l’extension d’Euler utilisent ce générateur après sa
construction modale.

La récurrence $(a_n)$ admet ensuite son identification à la
suite de Fibonacci. L’ordre démonstratif de cette sous-section est explicite :
sélecteur tritique, succession de modes, décompte des arêtes, incidence et
récurrence. L’identification de la suite conserve cet ordre. La
génération régionale coinductive possède en outre ses propres émetteurs, filtres,
relèvements et conditions de survie ; la récurrence modale précise
l’une de ses réalisations structurelles.

\subsection{Extension arithmétique du calendrier}

\subsubsection{Décomposition en phase et position dodécaphasique}

Nous écrivons $C_q=\ZZ/q\ZZ$ pour le groupe additif des résidus modulo $q$.
La position d’un calendrier de $108=12\cdot9$ événements possède une
décomposition unique
\[
 t=9b+r,\qquad 0\leq b<12,\quad0\leq r<9,
\]
pour le représentant $0\leq t<108$. La coordonnée $r$ est la phase
nonadique ; $b$ identifie la position du tour dans le registre
dodécaphasique. L’échelle $108$ réunit ces deux registres selon la loi de
composition avec retenue.

\begin{teorema}[Extension du calendrier avec retenue]
\label{thm:cal-extension}
Les applications
\[
 \iota:C_{12}\longrightarrow C_{108},\quad[b]\longmapsto[9b],
 \qquad p:C_{108}\longrightarrow C_9,\quad[t]\longmapsto[t]_9
\]
déterminent une suite exacte
\begin{equation}
 0\longrightarrow C_{12}\xrightarrow{\ \iota\ }C_{108}
 \xrightarrow{\ p\ }C_9\longrightarrow0.
 \label{eq:cal-extension}
\end{equation}
Dans les coordonnées $(b,r)$, l’addition est donnée par
\begin{equation}
 (b,r)\oplus(b',r')=
 \left(\left[b+b'+\left\lfloor\frac{r+r'}9\right\rfloor\right]_{12},
 [r+r']_9\right).
 \label{eq:cal-ley-grupo}
\end{equation}
Cette extension n’a pas de section homomorphe.
\end{teorema}
\begin{proof}
L’égalité $[9b]=0$ dans $C_{108}$ équivaut à $12\mid b$, donc
$\iota$ est injective. L’application $p$ est surjective et son noyau
est constitué exactement des classes multiples de neuf, image de
$\iota$. Cela démontre l’exactitude.

Pour additionner deux représentants, on écrit
\[
 9b+r+9b'+r'
 =9\left(b+b'+\left\lfloor\frac{r+r'}9\right\rfloor\right)
   +[r+r']_9.
\]
La réduction du premier coefficient modulo douze donne
\eqref{eq:cal-ley-grupo}. Supposons maintenant qu’il existe une section
homomorphe $s:C_9\to C_{108}$. Le représentant $a$ de $s([1])$
satisfait $a\equiv1\pmod9$ et $9a\equiv0\pmod{108}$. La seconde
condition implique $12\mid a$, d’où $a\equiv0\pmod3$ ; la première
implique $a\equiv1\pmod3$. La contradiction démontre la dernière
affirmation.
\end{proof}

\subsubsection{Cocycle, inverse et restitution}

Le terme
\[
 \kappa_9(r,s)=\left\lfloor\frac{r+s}9\right\rfloor,
 \qquad 0\leq r,s<9,
\]
enregistre le franchissement du seuil. Son identité de composition est
\begin{equation}
 \kappa_9(r,s)+\kappa_9([r+s]_9,u)
 =\kappa_9(s,u)+\kappa_9(r,[s+u]_9).
 \label{eq:cal-cociclo}
\end{equation}
Les deux membres sont $\lfloor(r+s+u)/9\rfloor$ : il suffit de substituer
$r+s=9\kappa_9(r,s)+[r+s]_9$, ou la décomposition analogue de $s+u$.
L’identité garantit que le quotient accumulé est indépendant du
parenthésage de trois termes.

L’inverse de $(b,r)$ est déterminée par
\[
 -(b,r)=
 \begin{cases}
 ([-b]_{12},0),&r=0,\\
 ([-b-1]_{12},9-r),&1\leq r<9.
 \end{cases}
\]
Le terme supplémentaire $-1$ représente la correction exacte de frontière lors de
l’inversion d’un résidu intérieur. Par exemple, le représentant $t=17$ a pour
coordonnées $(1,8)$. Son inverse est $(-2,1)$ dans les coordonnées dont la
première composante est modulo douze : $9(-2)+1=-17$. Le résidu inversé et
le quotient corrigé restituent conjointement l’entier opposé.

\subsection{Relèvement entier et holonomie nonadique}

\subsubsection{Phase finie et compteur de tours}

La division euclidienne s’étend à tout $k\in\ZZ$ :
\[
 k=9m+r,\qquad m=\left\lfloor\frac k9\right\rfloor,
 \quad0\leq r<9.
\]
La paire $(m,r)$ détermine exactement $k$. Sa réduction
$(m,r)\mapsto([m]_{12},r)$ retrouve le calendrier fini. Dans le
relèvement entier, un tour de neuf pas agit par
\begin{equation}
 (r,m)\longmapsto(r,m+1).
 \label{eq:cal-vuelta-entera}
\end{equation}
Douze tours restituent les deux coordonnées finies et translatent le
compteur entier de douze unités. La mémoire enrichie comprend ce
compteur ainsi que le registre ordonné des transitions.

\subsubsection{Composition des prolongations admissibles}

Au niveau $k$, nous représentons l’état régional enrichi par
\[
 \widehat x_k=(B_k,C_k,p_k,c_k,\Sigma_k,W_k,g_k,m_k).
\]
$B_k$ réunit les blocs régionaux ; $C_k$ est le cylindre compatible ;
$p_k$ est le préfixe ; $c_k$ enregistre la retenue ; $\Sigma_k$ conserve la
frontière et $W_k$ l’incidence. Les dernières coordonnées sont la phase et
le compteur de tours. Les opérateurs de prolongation agissent sur
cet état complet, avec les dépendances de chaque composante.

\begin{definicion}[Holonomie nonadique du système de prolongation]
Soient $\mathcal R_k\subseteq\widehat X_k\times\widehat X_{k+1}$
les relations admissibles produites par le TPK. La composition d’un
tour est
\[
 \Gamma_{9,k}=\mathcal R_{k+8}\circ\cdots\circ\mathcal R_k.
\]
Pour $s\geq1$, la composition de $s$ tours est
\[
 \Gamma_{9s,k}=
 \Gamma_{9,k+9(s-1)}\circ\cdots\circ\Gamma_{9,k}.
\]
\end{definicion}

La composition relationnelle retient les prolongations qui satisfont
simultanément les contraintes de l’état. Une trajectoire admissible
individualise une succession de telles transitions. La notation distingue
la relation d’admissibilité, la trajectoire choisie et l’observation
finie de sa phase.

\begin{proposicion}[Retour de phase et avancée de mémoire]
\label{prop:cal-memoria}
Sur une trajectoire dont la règle de tour est
$g_{k+9}=g_k$, $m_{k+9}=m_k+1$, toute composition admissible de $s$
tours satisfait
\[
 g_{k+9s}=g_k,\qquad m_{k+9s}=m_k+s,
 \qquad [m_{k+9s}]_{12}=[m_k+s]_{12}.
\]
\end{proposicion}
\begin{proof}
Le cas $s=1$ est la règle de tour. En composant un tour supplémentaire
avec la trajectoire déjà construite, la phase est conservée et le compteur
reçoit une unité de plus. La récurrence démontre les deux premières égalités ;
la projection modulo douze donne la troisième.
\end{proof}

Le registre permet de reconnaître chaque retour comme une transformation de
l’histoire. La phase observable exprime une périodicité, tandis que
le compteur et les coordonnées de frontière distinguent ses réalisations
successives. Les démonstrations de prolongation régionale précisent
en outre quels cylindres survivent et quels préfixes sont déterminés.

\subsection{Représentation bilatérale du transport de phase}

\subsubsection{Décalage unitaire et groupe des tours}

Le relèvement de l’indice admet une représentation dans
$\mathcal H_{\rm ind}=\ell^2(\ZZ)$. Sa base orthonormée s’écrit
$e_k=e_{r,m}$ lorsque $k=9m+r$. Nous définissons $Te_k=e_{k+1}$, c’est-à-dire
\[
 Te_{r,m}=\begin{cases}
 e_{r+1,m},&r<8,\\e_{0,m+1},&r=8.
 \end{cases}
\]
La transformation permute bijectivement une base orthonormée et se
prolonge par continuité en un opérateur unitaire. L’inverse envoie
$e_{r,m}$ sur $e_{r-1,m}$ lorsque $r>0$ et envoie $e_{0,m}$ sur
$e_{8,m-1}$.

\begin{teorema}[Représentation fidèle du compteur de tours]
\label{thm:cal-representacion}
Soit $S=T^9$. Alors
\[
 S^se_{r,m}=e_{r,m+s},\qquad s\in\ZZ.
\]
L’application $s\mapsto S^s$ est une représentation unitaire fidèle de
$(\ZZ,+)$.
\end{teorema}
\begin{proof}
Neuf pas conservent le résidu et augmentent le quotient d’une unité.
La même identité appliquée à l’inverse fournit la formule pour
les exposants négatifs. Les puissances satisfont
$S^{s+t}=S^sS^t$. Si $S^s=I$, l’égalité
$e_{r,m+s}=e_{r,m}$ impose $s=0$ par l’indépendance des vecteurs
de la base. Cela prouve la fidélité.
\end{proof}

L’espace des indices est une réalisation du compteur. L’espace des
états du TPK incorpore en outre les données arithmétiques et d’incidence.
L’inversion unitaire de l’indice s’applique à son revêtement bilatéral ;
l’inversion d’une transition complète utilise l’archive de cette
transition et ses conditions d’admissibilité.

\subsubsection{Opérateur de mémoire et correction de frontière}

Définissons, sur les vecteurs à support fini,
\[
 D_Me_k=\left\lfloor\frac k9\right\rfloor e_k,
 \qquad \mathcal Ie_k=e_{-k},
 \qquad P_0e_k=\begin{cases}e_k,&9\mid k,\\0,&9\nmid k.\end{cases}
\]
Ce domaine est dense et stable sous les opérations considérées.

\begin{proposicion}[Identités de mémoire et d’inversion]
Sur ce domaine, on a
\begin{equation}
 [D_M,S]=S,\qquad
 \mathcal IT\mathcal I=T^{-1},\qquad
 \mathcal ID_M\mathcal I=-D_M-(I-P_0).
 \label{eq:cal-inversion}
\end{equation}
\end{proposicion}
\begin{proof}
Si $k=9m+r$, alors
$D_MS e_k=(m+1)e_{k+9}$ et $SD_Me_k=me_{k+9}$, ce qui prouve
le commutateur. La composition $\mathcal IT\mathcal I$ envoie
$e_k$ sur $e_{k-1}$. Enfin,
\[
 \left\lfloor\frac{-k}9\right\rfloor
 =\begin{cases}-m,&r=0,\\-m-1,&r>0.\end{cases}
\]
L’action diagonale du dernier membre de \eqref{eq:cal-inversion}
coïncide avec cette expression sur chaque vecteur de la base. La linéarité
achève la démonstration.
\end{proof}

La correction $I-P_0$ situe la contribution de frontière dans les phases
intérieures. Lorsque le parcours est inversé, le quotient et le résidu se
transforment conjointement. Cette identité exprime arithmétiquement
l’information qui disparaît lorsque l’on n’observe que la phase.

\subsection{Capacité de raffinement et vacances arithmétiques}

\subsubsection{Comparaison exacte des blocs ternaires et décimaux}

Un bloc de six trits dispose de $3^6=729$ configurations dans son
espace ambiant. Un bloc décimal de trois chiffres dispose de $10^3=1000$
positions. Leurs cardinaux définissent une capacité arithmétique de
codage. Pour $n\geq0$, soit
\begin{equation}
 \mathcal C_n=\max\{a\in\NN:1000^a\leq729^{n+1}\}.
 \label{eq:cal-capacidad-entera}
\end{equation}
La définition compare des entiers et peut être évaluée exactement. Dans une
représentation logarithmique ultérieure,
\[
 \rho=\log_{1000}729=\log_{10}9,\qquad
 \delta=1-\rho,\qquad
 \mathcal C_n=\lfloor(n+1)\rho\rfloor.
\]
Le rapport $\rho$ appartient à $(0,1)$ et est irrationnel. Si
$\rho=p/q$ avec des entiers positifs, on aurait $10^p=9^q$ ;
la factorisation en nombres premiers des deux membres l’interdit.

\begin{definicion}[Registre des vacances de capacité]
On définit
\[
 Z_n=\lfloor(n+1)\delta\rfloor,\qquad
 z_n=\lfloor(n+1)\delta\rfloor-\lfloor n\delta\rfloor,
 \quad n\geq0.
\]
La valeur $z_n\in\{0,1\}$ enregistre la rétention de capacité au
pas correspondant. On a $z_0=0$ à l’origine choisie.
\end{definicion}

\begin{teorema}[Bilan de capacité et de vacances]
\label{thm:cal-capacidad}
Pour $n\geq0$ dans la première identité et $n\geq1$ dans la seconde,
\begin{equation}
 \mathcal C_n+Z_n=n,
 \qquad \mathcal C_n-\mathcal C_{n-1}=1-z_n.
 \label{eq:cal-capacidad}
\end{equation}
En particulier, pour $0\leq a<b$,
\[
 \mathcal C_b-\mathcal C_a+\sum_{n=a+1}^{b}z_n=b-a.
\]
\end{teorema}
\begin{proof}
L’irrationalité de $\delta$ implique
$\lceil(n+1)\delta\rceil=\lfloor(n+1)\delta\rfloor+1$.
Donc,
\[
 \lfloor(n+1)(1-\delta)\rfloor
 =n+1-\lceil(n+1)\delta\rceil=n-Z_n.
\]
La différence entre deux égalités consécutives donne la seconde identité
pour $n\geq1$. La somme télescopique démontre la formule sur un intervalle.
\end{proof}

Le registre $Z_n$ peut être obtenu au moyen de $n-\mathcal C_n$,
en utilisant uniquement la comparaison entière
\eqref{eq:cal-capacidad-entera}. Les fonctions logarithmiques expriment
la même construction en coordonnées analytiques. La première rétention
de capacité se produit entre $n=20$ et $n=21$ :
\[
 \mathcal C_{19}=19,\quad\mathcal C_{20}=20,\quad
 \mathcal C_{21}=20,\quad\mathcal C_{22}=21.
\]
L’histoire reçoit un bloc supplémentaire tandis que la capacité décimale
reste constante. La vacance enregistre précisément ce décalage.

\subsubsection{Différenciation des indices opératoires}

La phase de l’émetteur fini, la phase de prolongation régionale et la capacité
décimale appartiennent à des indices déclarés. La connexion nonadique conserve
la phase après neuf pas de son propre indice ; la vacance conserve une
capacité pendant un pas de l’indice des blocs. La comparaison entre
les deux s’exprime par
\[
 [\mathcal C_n]_q=\bigl[[n]_q-[Z_n]_q\bigr]_q,
 \qquad q=9\ \text{ou}\ 108.
\]
Le décalage de phase dépend du registre accumulé des vacances. Le conserver
permet de transformer une lecture de phase en l’autre sans identifier
indûment les deux calendriers.

La structure obtenue relie sélection modale, composition de la retenue,
holonomie et raffinement. La mémoire remplit des fonctions précises à chaque
niveau : elle restitue les quotients, distingue les tours, enregistre les transitions et
contrôle la capacité de codage. Ces opérations constituent le
support temporel et archivistique sur lequel agissent les lecteurs régionaux
et les réalisations de conservation de l’unité.
